\documentclass[a4paper]{arrowhead}

\usepackage[yyyymmdd]{datetime}
\usepackage{etoolbox}
\usepackage[T1]{fontenc}
\usepackage[utf8]{inputenc}
\usepackage{multirow}
\usepackage{longtable}
\usepackage{amssymb}
\usepackage{float}
\usepackage{tikz}
\usetikzlibrary{positioning,arrows.meta,fit,calc,shapes.misc,shapes.geometric}

\renewcommand{\dateseparator}{-}
\setlength{\emergencystretch}{3em}
\lstset{basicstyle=\ttfamily\footnotesize,breaklines=true,columns=fullflexible,frame=single,rulecolor=\color{gray!50},backgroundcolor=\color{gray!6}}
\tikzset{
  ahsys/.style={draw=ArrowheadBlue,very thick,fill=ArrowheadBlue!6,rounded corners=2pt,minimum width=3.2cm,minimum height=1.1cm,align=center,font=\small\bfseries},
  ahext/.style={draw=ArrowheadGrey,thick,dashed,fill=gray!5,rounded corners=2pt,minimum width=3.2cm,minimum height=1.0cm,align=center,font=\small},
  ahsrv/.style={draw=ArrowheadPurple,thick,fill=ArrowheadPurple!8,rounded corners=6pt,minimum height=0.55cm,inner xsep=4pt,align=center,font=\scriptsize\ttfamily},
  ahcsm/.style={draw=ArrowheadTeal,thick,fill=ArrowheadTeal!8,rounded corners=6pt,minimum height=0.55cm,inner xsep=4pt,align=center,font=\scriptsize\ttfamily},
  ahdb/.style={draw=ArrowheadOrange,thick,fill=ArrowheadOrange!8,cylinder,shape border rotate=90,aspect=0.25,minimum height=0.9cm,align=center,font=\scriptsize},
  ahflow/.style={-{Stealth[length=2.2mm]},thick,ArrowheadGrey},
  ahlabel/.style={font=\scriptsize,fill=white,inner sep=1pt,align=center},
}

\newcommand{\fparam}[1]{\textit{\textcolor{ArrowheadBlue}{#1}}}
%% Special references
\newcommand{\mref}[1]{{\textcolor{ArrowheadPurple}{\hyperref[sec:model:#1]{#1}}}}
\newcommand{\oref}[1]{{\textcolor{ArrowheadBlue}{\hyperref[sec:operations:#1]{#1}}}}
\newcommand{\pdef}[1]{{\textcolor{ArrowheadGrey}{#1\label{sec:model:primitives:#1}\label{sec:model:primitives:#1s}}}}
\newcommand{\pref}[1]{{\textcolor{ArrowheadGrey}{\hyperref[sec:model:primitives:#1]{#1}}}}

\newrobustcmd\osubsection[5]{
  \addtocounter{subsection}{1}
  \addcontentsline{toc}{subsection}{\protect\numberline{\thesubsection}operation \textcolor{ArrowheadBlue}{#1}}
  \subsection*{
    \thesubsection\quad
    #2 \textcolor{ArrowheadPurple}{#3} \\
    \small
    \hspace*{0.075\textwidth}\begin{minipage}{0.1\textwidth}
      \vspace*{1mm}
      Operation: \\
      \notblank{#4}{Input: \\}{}
      \notblank{#5}{Output: \\}{}
    \end{minipage}
    \begin{minipage}{0.825\textwidth}
      \vspace*{1mm}
      \textcolor{ArrowheadBlue}{#1} \\
      \notblank{#4}{\mref{#4} \\}{}
      \notblank{#5}{\mref{#5} \\}{}
    \end{minipage}
  }
  \label{sec:operations:#1}
}
\newrobustcmd\msubsection[2]{
  \addtocounter{subsection}{1}
  \addcontentsline{toc}{subsection}{\protect\numberline{\thesubsection}#1 \textcolor{ArrowheadPurple}{#2}}
  \subsection*{\thesubsection\quad#1 \textcolor{ArrowheadPurple}{#2}}
  \label{sec:model:#2} \label{sec:model:#2s}
}

\begin{document}

%% Arrowhead Document Properties
\ArrowheadTitle{workflowExecution HTTP/TLS/JSON}
\ArrowheadServiceID{workflowExecution}
\ArrowheadType{Interface Design Description}
\ArrowheadTypeShort{IDD}
\ArrowheadVersion{5.4.0}
\ArrowheadDate{2026-11-24}
\ArrowheadAuthor{Documentation maintenance}
\ArrowheadStatus{REVIEW DRAFT}
\ArrowheadContact{Repository maintainers}
\ArrowheadFooter{\href{www.arrowhead.eu}{www.arrowhead.eu}}
\ArrowheadSetup
%%

%% Front Page
\begin{center}
  \vspace*{1cm}
  \huge{\arrowtitle}

  \vspace*{0.2cm}
  \LARGE{\arrowtype}
  \vspace*{1cm}

  \Large{Service ID: \textit{``\arrowid''}}
  \vspace*{\fill}

  

  \vspace*{1cm}
  \vspace*{\fill}

  % Front Page Abstract
  \begin{abstract}
    This document describes the HTTP/TLS/JSON interface of the \texttt{workflow\allowbreak{}Execution} service of the Eclipse Arrowhead workflow support systems, service version 1.0.0, documentation edition 5.4.0. It realises the abstract operations of the \texttt{workflow\allowbreak{}Execution} SD.
  \end{abstract}

  \vspace*{1cm}
\end{center}

\newpage
%%

%% Table of Contents
\tableofcontents
\newpage
%%

\section{Overview}
\label{sec:overview}
This document describes how the \texttt{workflow\allowbreak{}Execution} service, which runs work items on behalf of a WorkflowManager, is realised with HTTP~\cite{rfc9110}
over TLS~\cite{rfc8446} and JSON~\cite{rfc8259} payloads. Request and response bodies use content type
\texttt{application/\allowbreak{}json}; JSON field names are case-sensitive and fields not listed in this document must not be sent
(the recipe decoder rejects unknown fields). Optional fields may be omitted; absent and \texttt{null} are equivalent.

The service is offered by the WorkflowExecutor on its Arrowhead connector (default port 8483) with base path \texttt{/\allowbreak{}workflowexecutor/\allowbreak{}execution}.
The ServiceRegistry entry lists every operation with its HTTP method and path relative to the base path.

The rest of this document is organized as follows. Section~\ref{sec:operations} specifies each operation,
Section~\ref{sec:model} the JSON data models, Section~\ref{sec:security} the security requirements,
Section~\ref{sec:errors} the error representation and Section~\ref{sec:implementations} the implementations.

\newpage
\section{Service Operations}
\label{sec:operations}

{\small
\begin{longtable}{|p{3.0cm}|p{1.5cm}|p{5.3cm}|p{3.0cm}|p{3.0cm}|}

\hline
\rowcolor{gray!33} Operation & Method & Path & Input & Output \\ \hline
\endhead
\oref{service-info} & POST & \texttt{/\allowbreak{}workflowexecutor/\allowbreak{}execution/\allowbreak{}service-\allowbreak{}info} & -- & \mref{ExecutorInfo} \\ \hline
\oref{start} & POST & \texttt{/\allowbreak{}workflowexecutor/\allowbreak{}execution/\allowbreak{}start} & \mref{WorkItem} & -- \\ \hline
\oref{abort} & POST & \texttt{/\allowbreak{}workflowexecutor/\allowbreak{}execution/\allowbreak{}abort} & \mref{AbortRequest} & -- \\ \hline
\oref{get-status} & GET & \texttt{/\allowbreak{}workflowexecutor/\allowbreak{}execution/\allowbreak{}status/\allowbreak{}\{workItemId\}} & \mref{WorkItemRef} & \mref{ExecutionState} \\ \hline
\end{longtable}
}

\osubsection{service-info}{POST}{/workflowexecutor/execution/service-info}{}{ExecutorInfo}
Returns the executor's name, capabilities $cap(w)$, runnable services, workstation and mode.

{\small
\begin{longtable}{|p{2.0cm}|p{13.8cm}|}

\hline
\rowcolor{gray!33} Status & When \\ \hline
\endhead
Success & 200 OK. \\ \hline
401 & No valid identity token, or the token is bound to another certificate (\texttt{AUTH}). \\ \hline
403 & The consumer has no ConsumerAuthorization grant for this operation (\texttt{FORBIDDEN}). \\ \hline
\end{longtable}
}


\begin{lstlisting}[caption={Example response body}]
{ "name": "WfeA", "capabilities": ["heating"], "services": ["*"],
  "workstation": "cell-1", "mode": "push" }
\end{lstlisting}

\osubsection{start}{POST}{/workflowexecutor/execution/start}{WorkItem}{}
Accepts a work item and runs it asynchronously. The outcome is reported through \texttt{workflow\allowbreak{}Work\allowbreak{}Item.\allowbreak{}report}.

{\small
\begin{longtable}{|p{2.0cm}|p{13.8cm}|}

\hline
\rowcolor{gray!33} Status & When \\ \hline
\endhead
Success & 202 Accepted, empty body. \\ \hline
400 & Work item incomplete (\texttt{workItemId} or \texttt{step} missing). \\ \hline
403 & The caller is not the configured WorkflowManager. \\ \hline
401 & No valid identity token, or the token is bound to another certificate (\texttt{AUTH}). \\ \hline
403 & The consumer has no ConsumerAuthorization grant for this operation (\texttt{FORBIDDEN}). \\ \hline
\end{longtable}
}

\osubsection{abort}{POST}{/workflowexecutor/execution/abort}{AbortRequest}{}
Aborts a running work item; the executor reports ABORTED.

{\small
\begin{longtable}{|p{2.0cm}|p{13.8cm}|}

\hline
\rowcolor{gray!33} Status & When \\ \hline
\endhead
Success & 200 OK, empty body. \\ \hline
403 & The caller is not the configured WorkflowManager. \\ \hline
401 & No valid identity token, or the token is bound to another certificate (\texttt{AUTH}). \\ \hline
403 & The consumer has no ConsumerAuthorization grant for this operation (\texttt{FORBIDDEN}). \\ \hline
\end{longtable}
}


\begin{lstlisting}[caption={Example request body}]
{ "workItemId": "5d2c0b1e-...", "sessionId": 42, "reason": "session aborted" }
\end{lstlisting}

\osubsection{get-status}{GET}{/workflowexecutor/execution/status/\{workItemId\}}{WorkItemRef}{ExecutionState}
Returns the executor's local state of a work item. The WorkflowManager uses it as liveness check: a held work item is lost when two checks in a row find it neither \texttt{ACCEPTED} nor \texttt{RUNNING}, or when the executor cannot be reached or is no longer registered for \texttt{workflow.\allowbreak{}executor.\allowbreak{}liveness.\allowbreak{}lost-\allowbreak{}after-\allowbreak{}seconds}. A lost item that is not bound goes to the next executor; a bound one is handled like \texttt{ERROR}.

{\small
\begin{longtable}{|p{3.0cm}|p{1.4cm}|p{3.2cm}|p{8.2cm}|}

\hline
\rowcolor{gray!33} Parameter & In & Type & Description \\ \hline
\endhead
\fparam{workItemId} & path & \pref{WorkItemId} & Work item. \\ \hline
\end{longtable}
}


{\small
\begin{longtable}{|p{2.0cm}|p{13.8cm}|}

\hline
\rowcolor{gray!33} Status & When \\ \hline
\endhead
Success & 200 OK. \\ \hline
401 & No valid identity token, or the token is bound to another certificate (\texttt{AUTH}). \\ \hline
403 & The consumer has no ConsumerAuthorization grant for this operation (\texttt{FORBIDDEN}). \\ \hline
\end{longtable}
}


\begin{lstlisting}[caption={Example response body}]
{ "workItemId": "5d2c0b1e-...", "state": "RUNNING" }
\end{lstlisting}


\newpage
\section{Data Models}
\label{sec:model}
The JSON objects exchanged by the operations, in alphabetic order. Field types refer to other data models or to the
primitives of Section~\ref{sec:model:primitives}.

\msubsection{struct}{AbortRequest}
Aborts a work item on an executor.

{\small
\begin{longtable}{|p{3.7cm}|p{4.2cm}|p{1.6cm}|p{6.3cm}|}

\hline
\rowcolor{gray!33} Field & Type & Mandatory & Description \\ \hline
\endhead
\texttt{workItemId} & \pref{WorkItemId} & yes & The work item. \\ \hline
\texttt{sessionId} & \pref{SessionId} & yes & Its session. \\ \hline
\texttt{reason} & \pref{String} & no & Why. \\ \hline
\end{longtable}
}

\msubsection{enum}{DispatchMode}
How a step reaches an executor.
 Encoded as a JSON string.
{\small
\begin{longtable}{|p{3.9cm}|p{11.9cm}|}

\hline
\rowcolor{gray!33} Value & Meaning \\ \hline
\endhead
\texttt{push} & The WorkflowManager selects an executor and calls \texttt{workflow\allowbreak{}Execution.\allowbreak{}start} (default). \\ \hline
\texttt{pull} & The WorkflowManager publishes the work item; an executor claims it with \texttt{workflow\allowbreak{}Work\allowbreak{}Item.\allowbreak{}claim}. \\ \hline
\texttt{event} & Reserved for EventHandler notification; currently handled as \texttt{pull}. \\ \hline
\end{longtable}
}

\msubsection{struct}{ErrorMessage}
Error response of every operation (common library \texttt{Error\allowbreak{}MessageDTO}).

{\small
\begin{longtable}{|p{3.7cm}|p{4.2cm}|p{1.6cm}|p{6.3cm}|}

\hline
\rowcolor{gray!33} Field & Type & Mandatory & Description \\ \hline
\endhead
\texttt{errorMessage} & \pref{String} & yes & Explanation. \\ \hline
\texttt{errorCode} & \pref{Integer} & yes & HTTP status code. \\ \hline
\texttt{exceptionType} & \pref{String} & yes & \texttt{INVALID\_\allowbreak{}PARAMETER}, \texttt{AUTH}, \texttt{FORBIDDEN}, \texttt{DATA\_\allowbreak{}NOT\_\allowbreak{}FOUND}, \texttt{EXTERNAL\_\allowbreak{}SERVER\_\allowbreak{}ERROR}, \ldots \\ \hline
\texttt{origin} & \pref{String} & no & The operation that failed, e.g. \texttt{POST /\allowbreak{}workflowmanager/\allowbreak{}work-\allowbreak{}item/\allowbreak{}report}. \\ \hline
\end{longtable}
}

\msubsection{struct}{ExecutionState}
State of a work item on an executor.

{\small
\begin{longtable}{|p{3.7cm}|p{4.2cm}|p{1.6cm}|p{6.3cm}|}

\hline
\rowcolor{gray!33} Field & Type & Mandatory & Description \\ \hline
\endhead
\texttt{workItemId} & \pref{WorkItemId} & yes & The work item. \\ \hline
\texttt{state} & \mref{ExecutorState} & yes & Local state. \\ \hline
\texttt{message} & \pref{String} & no & Last message. \\ \hline
\end{longtable}
}

\msubsection{struct}{ExecutorInfo}
What an executor declares about itself.

{\small
\begin{longtable}{|p{3.7cm}|p{4.2cm}|p{1.6cm}|p{6.3cm}|}

\hline
\rowcolor{gray!33} Field & Type & Mandatory & Description \\ \hline
\endhead
\texttt{name} & \pref{SystemName} & yes & Executor system name. \\ \hline
\texttt{capabilities} & \pref{List}$<$\pref{String}$>$ & yes & Capabilities $cap(w)$. \\ \hline
\texttt{services} & \pref{List}$<$\pref{ServiceDefinition}$>$ & yes & Services it can run; \texttt{*} means any. \\ \hline
\texttt{workstation} & \pref{String} & no & Workstation. \\ \hline
\texttt{mode} & \mref{DispatchMode} & yes & \texttt{push} or \texttt{pull}. \\ \hline
\end{longtable}
}

\msubsection{struct}{ExecutorRequirement}
$\eta(s)$: constraints on the executor that runs a step.

{\small
\begin{longtable}{|p{3.7cm}|p{4.2cm}|p{1.6cm}|p{6.3cm}|}

\hline
\rowcolor{gray!33} Field & Type & Mandatory & Description \\ \hline
\endhead
\texttt{capability} & \pref{String} & no & A capability the executor must declare in \texttt{workflow\allowbreak{}Capabilities}. \\ \hline
\texttt{workstation} & \pref{String} & no & The executor's workstation. \\ \hline
\texttt{system} & \pref{SystemName} & no & A named executor system. \\ \hline
\end{longtable}
}

\msubsection{enum}{ExecutorState}
Local state of a work item on an executor. The authoritative state is held by the WorkflowManager.
 Encoded as a JSON string.
{\small
\begin{longtable}{|p{3.9cm}|p{11.9cm}|}

\hline
\rowcolor{gray!33} Value & Meaning \\ \hline
\endhead
\texttt{ACCEPTED} & Received, not yet started. \\ \hline
\texttt{RUNNING} & Being run. \\ \hline
\texttt{DONE} & Run finished with \texttt{SUCCESS}. \\ \hline
\texttt{BIND\_FAILED} & Run ended with this report. \\ \hline
\texttt{ERROR} & Run ended with this report. \\ \hline
\texttt{FATAL\_ERROR} & Run ended with this report. \\ \hline
\texttt{ABORTED} & Run aborted. \\ \hline
\texttt{UNKNOWN} & The executor does not know the work item, e.g. after a restart. \\ \hline
\end{longtable}
}

\msubsection{enum}{LeaseMode}
Role of a step for a shared resource.
 Encoded as a JSON string.
{\small
\begin{longtable}{|p{3.9cm}|p{11.9cm}|}

\hline
\rowcolor{gray!33} Value & Meaning \\ \hline
\endhead
\texttt{reserve} & Reserves the resource; its result carries the lease token. \\ \hline
\texttt{use} & Uses the resource; the executor renews the lease immediately before the invocation. \\ \hline
\texttt{release} & Releases the resource. \\ \hline
\end{longtable}
}

\msubsection{struct}{LeasePolicy}
The role of a step for a shared resource.

{\small
\begin{longtable}{|p{3.7cm}|p{4.2cm}|p{1.6cm}|p{6.3cm}|}

\hline
\rowcolor{gray!33} Field & Type & Mandatory & Description \\ \hline
\endhead
\texttt{resource} & \pref{Name} & yes & A declared resource. \\ \hline
\texttt{mode} & \mref{LeaseMode} & yes & \texttt{reserve}, \texttt{use} or \texttt{release}. \\ \hline
\end{longtable}
}

\msubsection{enum}{OnError}
Reaction to an \texttt{ERROR} report.
 Encoded as a JSON string.
{\small
\begin{longtable}{|p{3.9cm}|p{11.9cm}|}

\hline
\rowcolor{gray!33} Value & Meaning \\ \hline
\endhead
\texttt{rebind} & Exclude the executor for this session and bind the step again. \\ \hline
\texttt{retry} & Keep the executor: after \texttt{backoff\allowbreak{}Seconds} dispatch the step to the same executor again, or to the least-loaded candidate when that executor is no longer available (extension x-retry; for transient provider-side failures, the provider must accept a repeated call). \\ \hline
\texttt{abort} & Fail the session. \\ \hline
\end{longtable}
}

\msubsection{struct}{Policy}
Execution policy of a step. Every field other than \texttt{lease} and \texttt{local\allowbreak{}Cloud\allowbreak{}Only} is an extension outside the proven model.

{\small
\begin{longtable}{|p{3.7cm}|p{4.2cm}|p{1.6cm}|p{6.3cm}|}

\hline
\rowcolor{gray!33} Field & Type & Mandatory & Description \\ \hline
\endhead
\texttt{quantity} & \pref{Integer} & no & Number of executions; only 1 is accepted (check M01). \\ \hline
\texttt{timeout\allowbreak{}Seconds} & \pref{Duration} & no & Extension x-timeout: the step fails when it runs longer. \\ \hline
\texttt{nominal\allowbreak{}Duration\allowbreak{}Seconds} & \pref{Duration} & no & Expected duration, used to check lease durations (check L04). \\ \hline
\texttt{retry} & \mref{Retry} & no & Extensions x-rebind and x-retry: behaviour on \texttt{ERROR}. \\ \hline
\texttt{lease} & \mref{LeasePolicy} & no & The step's role for a shared resource. \\ \hline
\texttt{local\allowbreak{}Cloud\allowbreak{}Only} & \pref{Boolean} & no & Whether orchestration may return providers of other clouds (default true). \\ \hline
\end{longtable}
}

\msubsection{struct}{Retry}
Extensions x-rebind and x-retry.

{\small
\begin{longtable}{|p{3.7cm}|p{4.2cm}|p{1.6cm}|p{6.3cm}|}

\hline
\rowcolor{gray!33} Field & Type & Mandatory & Description \\ \hline
\endhead
\texttt{maxAttempts} & \pref{Integer} & no & Maximum attempts of the step in one session, counting every work item of the step; at least 1, default 1 (no further attempt). \\ \hline
\texttt{onError} & \mref{OnError} & no & What an \texttt{ERROR} report causes while attempts remain.; default \texttt{rebind}. \\ \hline
\texttt{backoff\allowbreak{}Seconds} & \pref{Duration} & no & Pause before a same-executor retry (\texttt{onError} = \texttt{retry}); 0--600, default 1. \\ \hline
\end{longtable}
}

\msubsection{struct}{ServiceRequirement}
Requirements on the provider of the step's service, passed to orchestration.

{\small
\begin{longtable}{|p{3.7cm}|p{4.2cm}|p{1.6cm}|p{6.3cm}|}

\hline
\rowcolor{gray!33} Field & Type & Mandatory & Description \\ \hline
\endhead
\texttt{service\allowbreak{}Definition} & \pref{ServiceDefinition} & yes & The service to consume. \\ \hline
\texttt{interfaces} & \pref{List}$<$\pref{String}$>$ & no & Accepted interface template names, e.g. \texttt{generic\_\allowbreak{}https}. \\ \hline
\texttt{minVersion} & \pref{String} & no & Lowest accepted service version. \\ \hline
\texttt{maxVersion} & \pref{String} & no & Highest accepted service version. \\ \hline
\texttt{metadata} & \pref{Map}$<$\pref{String}, \pref{JsonValue}$>$ & no & Metadata requirements on the service instance. \\ \hline
\texttt{qos} & \pref{Map}$<$\pref{String}, \pref{JsonValue}$>$ & no & QoS requirements forwarded to orchestration. \\ \hline
\texttt{requires\allowbreak{}Results} & \pref{List}$<$\pref{Name}$>$ & no & Steps whose results the provider needs; each must be in $dep(s)$ (check D04). \\ \hline
\end{longtable}
}

\msubsection{struct}{WorkItem}
One attempt to run one step: what an executor needs to run it.

{\small
\begin{longtable}{|p{3.7cm}|p{4.2cm}|p{1.6cm}|p{6.3cm}|}

\hline
\rowcolor{gray!33} Field & Type & Mandatory & Description \\ \hline
\endhead
\texttt{workItemId} & \pref{WorkItemId} & yes & Identifier of the attempt. \\ \hline
\texttt{sessionId} & \pref{SessionId} & yes & Session. \\ \hline
\texttt{recipe} & \pref{String} & yes & \texttt{name@version} of the recipe. \\ \hline
\texttt{action} & \pref{Name} & yes & Action of the step. \\ \hline
\texttt{step} & \pref{Name} & yes & Step name. \\ \hline
\texttt{attempt} & \pref{Integer} & yes & Attempt number from 1. \\ \hline
\texttt{service} & \mref{ServiceRequirement} & yes & $\theta(s)$. \\ \hline
\texttt{executor} & \mref{ExecutorRequirement} & no & $\eta(s)$. \\ \hline
\texttt{params} & \pref{Map}$<$\pref{String}, \pref{JsonValue}$>$ & no & Static parameters of the step. \\ \hline
\texttt{input} & \pref{Template} & no & Input template $\kappa_s$. \\ \hline
\texttt{policy} & \mref{Policy} & no & Step policy. \\ \hline
\texttt{lease\allowbreak{}Duration\allowbreak{}Seconds} & \pref{Duration} & no & Lease duration of the step's resource. \\ \hline
\texttt{parameters} & \pref{JsonObject} & no & Session parameters $\rho$ (included for convenience; also available from \texttt{workflow\allowbreak{}Context}). \\ \hline
\texttt{dependency\allowbreak{}Results} & \pref{Map}$<$\pref{Name}, \pref{JsonValue}$>$ & no & $\xi|dep(s)$ (also available from \texttt{workflow\allowbreak{}Context}). \\ \hline
\end{longtable}
}

\msubsection{struct}{WorkItemRef}
Identifies a work item.

{\small
\begin{longtable}{|p{3.7cm}|p{4.2cm}|p{1.6cm}|p{6.3cm}|}

\hline
\rowcolor{gray!33} Field & Type & Mandatory & Description \\ \hline
\endhead
\texttt{workItemId} & \pref{WorkItemId} & yes & The work item. \\ \hline
\end{longtable}
}


\subsection{Primitives}
\label{sec:model:primitives}
The abstract primitive types of the SD are represented in JSON~\cite{rfc8259} as follows.
\subsubsection*{alias \pdef{Boolean}}
JSON \texttt{true} or \texttt{false}.

\subsubsection*{alias \pdef{Duration}}
JSON number of seconds.

\subsubsection*{alias \pdef{Integer}}
JSON number without fraction.

\subsubsection*{alias \pdef{JsonObject}}
JSON object.

\subsubsection*{alias \pdef{JsonValue}}
Any JSON value~\cite{rfc8259}.

\subsubsection*{alias \pdef{List}}
JSON array.

\subsubsection*{alias \pdef{Map}}
JSON object whose member names are the keys.

\subsubsection*{alias \pdef{Name}}
JSON string matching \texttt{\^{}[A-Za-z][A-Za-z0-9\_-]\{0,62\}\$}.

\subsubsection*{alias \pdef{ServiceDefinition}}
JSON string.

\subsubsection*{alias \pdef{SessionId}}
JSON number (64-bit integer); a path segment when used in a path.

\subsubsection*{alias \pdef{String}}
JSON string (UTF-8).

\subsubsection*{alias \pdef{SystemName}}
JSON string.

\subsubsection*{alias \pdef{Template}}
Any JSON value; references are JSON strings.

\subsubsection*{alias \pdef{WorkItemId}}
JSON string holding a UUID; a path segment when used in a path.



\section{Security}
\label{sec:security}
The interface is offered only over TLS~\cite{rfc8446} (TLS 1.3 or 1.2) with mutual authentication: the consumer must
present an Arrowhead certificate issued under the local cloud's certificate authority, and certificates on the
cloud's CRL are refused. The service is registered with interface template \texttt{generic\_\allowbreak{}https} and interface
policy \texttt{CERT\_AUTH}.

Every request must carry the consumer's identity token, obtained from the Authentication system and bound to the
consumer's certificate:
\begin{lstlisting}
Authorization: Bearer IDENTITY-TOKEN//<identity token>
\end{lstlisting}
A token issued to another system's certificate is refused (401). The provider then checks a ConsumerAuthorization
grant for the operation (\texttt{authorization.\allowbreak{}verify}, \texttt{service.\allowbreak{}authorization.\allowbreak{}enabled=\allowbreak{}true}); a missing grant is
refused with 403. Each operation is authorized separately; \texttt{start} and \texttt{abort} additionally require the caller to be the configured WorkflowManager or a sysop.


\section{Errors}
\label{sec:errors}
Every error response carries an \mref{ErrorMessage} body with content type \texttt{application/\allowbreak{}json}, for example:
\begin{lstlisting}
{ "errorMessage": "Work item 5d2c0b1e-... is held by another executor",
  "errorCode": 403, "exceptionType": "FORBIDDEN",
  "origin": "POST /workflowmanager/work-item/report" }
\end{lstlisting}
Status codes follow HTTP semantics~\cite{rfc9110}: 400 invalid parameter, 401 authentication, 403 forbidden, 404 not
found, 409 conflict with the current state, 500 internal error and 503 failure of an external system.


\section{Implementations}
\label{sec:implementations}
\begin{itemize}
\item \textbf{Java (reference).} \texttt{Eclipse-\allowbreak{}Arrowhead-\allowbreak{}java/\allowbreak{}ah5-\allowbreak{}Workflow-\allowbreak{}java}, module \texttt{workflow-\allowbreak{}executor}. Paths, methods and service
operations are defined once in \texttt{workflow-\allowbreak{}model} (\texttt{wire/\allowbreak{}Workflow\allowbreak{}Services.\allowbreak{}java}, \texttt{wire/\allowbreak{}Service\allowbreak{}Operations.\allowbreak{}java}) and checked by the provider-side authorization filter.
\item \textbf{Go.} Not implemented: \texttt{Eclipse-\allowbreak{}Arrowhead-\allowbreak{}go} has no workflow systems. A Go implementation must
provide exactly the paths, methods, JSON field names and status codes of this IDD to be interoperable.
\item \textbf{Rust.} Not implemented: \texttt{Eclipse-\allowbreak{}Arrowhead-\allowbreak{}rust} has no workflow systems. The same requirement applies.
\end{itemize}
Until the Go and Rust implementations exist, the workflow systems are \emph{not} interoperable across the three
implementations; only the Java systems interoperate with Java, Go and Rust core systems through the common Arrowhead
interfaces (ServiceRegistry, Authentication, ConsumerAuthorization, DynamicServiceOrchestration).


\newpage
\bibliographystyle{IEEEtran}
\bibliography{bibliography}

\newpage

\section{Revision History}
\subsection{Amendments}

\noindent\begin{tabularx}{\textwidth}{| p{1cm} | p{2.2cm} | p{1.6cm} | X | p{3.6cm} |} \hline
\rowcolor{gray!33} No. & Date & Version & Subject of Amendments & Author \\ \hline
1 & 2026-11-23 & \arrowversion & Initial edition: interface design description of workflowExecution over HTTP/TLS/JSON, written from the Java reference implementation (ah5-Workflow-java). & Documentation maintenance \\ \hline
2 & 2026-11-24 & \arrowversion & Retry policy; executor states; get-status as liveness check. & Documentation maintenance \\ \hline
\end{tabularx}

\subsection{Quality Assurance}

\noindent\begin{tabularx}{\textwidth}{| p{1cm} | p{2.2cm} | p{1.6cm} | X |} \hline
\rowcolor{gray!33} No. & Date & Version & Approved by \\ \hline
1 & & \arrowversion & Pending review \\ \hline
\end{tabularx}

\end{document}
